% GENERATED FILE. Edit canonical lessons, then run npm run generate:books.
% canonical-source-hash: fnv1a64:a4241364165e97e6
% canonical-lessons: GE-C258-ich-war-froh, GE-C258-ich-hoffe-dass, GE-C258-review

\chapter{How It Felt — ich war froh, ich hoffe}
\label{ch:ge-how-it-felt}
\hlchaptermodality{258}

\begin{chapteropening}
\textbf{By the end of this chapter:} \emph{I can say how I felt with ich war and a reason, and say what I hope for and dream of.}

Complete GE-C258-review: how you felt on a journey and why, with dass or weil, and one hope or dream with zu.
\end{chapteropening}

\section[ich war froh, dass …]{ich war froh — how it felt}

\label{lesson:GE-C258-ich-war-froh}



\begin{quote}
Say \textquotedblleft{}I was tired\textquotedblright{}. (\emph{Ich war müde.}) That \emph{war} is a Präteritum German says out loud, in speech as well as in books.
\end{quote}

\subsection*{You'll want to know: froh}
\textbf{froh} — \textquotedblleft{}glad\textquotedblright{}: \emph{froh}, with a long \emph{o}.

\begin{quote}
\emph{Ich war froh.} — I was glad.
\end{quote}

\emph{Froh} and \emph{glücklich} are close. \emph{Glücklich} is happy, a state that can fill a day or a life. \emph{Froh} is often glad about one thing, and relieved: the train came after all, the exam is over.

\begin{grammarlens}[title={Grammar Lens: the feeling, then the reason}]
\emph{Ich war} takes any feeling you own: \emph{froh}, \emph{glücklich}, \emph{traurig}, \emph{stolz}, \emph{nervös}, \emph{wütend}, \emph{müde}. Then the reason, with \textbf{dass} or \textbf{weil}, and the verb at the end:

\begin{quote}
\emph{Ich war froh, dass der Zug endlich} \textbf{kam}. — I was glad the train came at last. \emph{Ich war traurig, weil meine Großmutter krank} \textbf{war}. — I was sad because my grandmother was ill. \emph{Ich war nervös, weil der Zug Verspätung} \textbf{hatte}. — I was nervous because the train was late.
\end{quote}

\emph{Dass} gives the thing you feel about; \emph{weil} gives the cause. After \emph{froh}, \emph{dass} is the usual choice. In the present it works the same way: \emph{Ich bin froh, dass du kommst.}
\end{grammarlens}

\begin{cousinweb}
\textbf{froh} $\leftarrow$ Middle High German \emph{vrō} $\leftarrow$ Old High German \emph{frō} (EtymDB 2.1). You met it inside \emph{freut mich} in chapter two. EtymDB also sets English \emph{frolic} beside it and links \emph{frolic} to Dutch \emph{vrolijk}, \textquotedblleft{}cheerful\textquotedblright{}; the usual account is that English borrowed the word from Dutch.
\end{cousinweb}

\subsection*{Your turn}
Say these aloud:

\begin{itemize}
\raggedright
  \item \emph{froh}
  \item \emph{Ich war froh, dass der Zug endlich kam.}
  \item \emph{Ich war traurig, weil meine Großmutter krank war.}
  \item \emph{Ich war stolz. Ich war müde, aber glücklich.}
  \item one feeling of your own from last week, with \emph{weil}
\end{itemize}

\begin{tcolorbox}[breakable,colback=teal!4,colframe=teal!35!black,title={Before you move on}]
\textquotedblleft{}I was glad\textquotedblright{}? (\emph{\textbf{Ich war froh.}}) Glad about one thing, or happy through and through: \emph{froh} or \emph{glücklich}? (\emph{\textbf{Froh}} for the one thing.) After \emph{dass} and \emph{weil}, where is the verb? (\textbf{At the end.})
\end{tcolorbox}
\section[ich hoffe, dass …; mein Traum ist es …]{ich hoffe, dass — hopes and dreams}

\label{lesson:GE-C258-ich-hoffe-dass}



\begin{quote}
Say \textquotedblleft{}I was glad the train came at last\textquotedblright{}. (\emph{Ich war froh, dass der Zug endlich kam.}) That looks back. This lesson looks forward.
\end{quote}

\begin{grammarlens}[title={Grammar Lens: ich hoffe, two ways}]
\textbf{hoffen}, to hope, takes a \emph{dass} clause like \emph{finden} and \emph{glauben}:

\begin{quote}
\emph{Ich hoffe, dass du bald} \textbf{kommst}. — I hope you come soon.
\end{quote}

When the one hoping and the one doing are the same person, German can drop \emph{dass} and use \textbf{zu} with the infinitive at the end:

\begin{quote}
\emph{Ich hoffe, dich bald} \textbf{zu besuchen}. — I hope to visit you soon.
\end{quote}

That \emph{zu} is the German \textquotedblleft{}to\textquotedblright{} of \textquotedblleft{}to visit\textquotedblright{}. It stands right before the infinitive, and both go to the end.
\end{grammarlens}

\begin{culture}
For a bigger ambition, German builds a frame around the dream:

\begin{quote}
\emph{Mein Traum ist es, nach Japan} \textbf{zu reisen}. — My dream is to travel to Japan.
\end{quote}

The \emph{es} holds a place for the dream, and the comma opens it: \emph{nach Japan zu reisen}. The infinitive goes last, with \emph{zu} in front of it.

With a separable verb, \emph{zu} slips inside, between prefix and verb: \emph{mitnehmen}, \textbf{mitzunehmen}. \emph{Ich hoffe, meine Familie mitzunehmen.}

You can say the same dream more plainly, too: \emph{Ich möchte nach Japan reisen.} The frame with \emph{es} puts more weight on the dream.
\end{culture}

\begin{cousinweb}
\textbf{Traum} $\leftarrow$ Middle High German \emph{troum}, from Proto-Germanic \emph{*draumaz} (EtymDB 2.1): the same word as English \emph{dream}. \textbf{hoffen}, from chapter one hundred and thirty-six, came into Middle High German from Low German, beside English \emph{hope} and Dutch \emph{hopen}.
\end{cousinweb}

\subsection*{Your turn}
Say these aloud:

\begin{itemize}
\raggedright
  \item \emph{Ich hoffe, dass du bald kommst.}
  \item \emph{Ich hoffe, dich bald zu besuchen.}
  \item \emph{Mein Traum ist es, nach Japan zu reisen.}
  \item \emph{mitnehmen}, then \emph{mitzunehmen}
  \item your own dream: \emph{Mein Traum ist es, … zu …}
\end{itemize}

\begin{tcolorbox}[breakable,colback=teal!4,colframe=teal!35!black,title={Before you move on}]
\textquotedblleft{}I hope to visit you soon\textquotedblright{}? (\emph{\textbf{Ich hoffe, dich bald zu besuchen.}}) Where does \emph{zu} go with \emph{mitnehmen}? (\textbf{Inside}: \emph{\textbf{mitzunehmen}}.) \textquotedblleft{}My dream is to travel to Japan\textquotedblright{}? (\emph{\textbf{Mein Traum ist es, nach Japan zu reisen.}})
\end{tcolorbox}
\section[(review)]{Review — a feeling, and a hope}

\label{lesson:GE-C258-review}



\begin{quote}
Say \textquotedblleft{}I was glad\textquotedblright{}. (\emph{Ich war froh.}) And \textquotedblleft{}I was tired, but happy\textquotedblright{}. (\emph{Ich war müde, aber glücklich.})
\end{quote}

\begin{grammarlens}[title={What you've built}]
Two lessons, and an experience can now have an inside. \emph{Ich war} and a feeling, then the reason with \emph{dass} or \emph{weil}. And a hope, with \emph{ich hoffe, dass …} or \emph{ich hoffe, … zu …}, or a dream framed by \emph{mein Traum ist es}.

The end of a story can carry both:

\begin{quote}
\emph{Schließlich war es Abend. Ich war froh, dass ich in Berlin war. Ich hoffe, meine Großmutter bald wieder zu besuchen.}
\end{quote}
\end{grammarlens}

\subsection*{Your turn}
Say these aloud:

\begin{itemize}
\raggedright
  \item \emph{Ich war froh, dass der Zug endlich kam.}
  \item \emph{Ich war traurig, weil es regnete.}
  \item \emph{Ich hoffe, dass du bald kommst.}
  \item \emph{Mein Traum ist es, nach Japan zu reisen.}
  \item the end of the journey from chapter two hundred and fifty-three, with one feeling and one hope
  \item a problem first, then the feeling: \emph{Ich habe den Zug verpasst. Was kann ich machen? Ich war sehr nervös.}
\end{itemize}

\begin{tcolorbox}[breakable,colback=teal!4,colframe=teal!35!black,title={Before you move on}]
\textquotedblleft{}I was glad that …\textquotedblright{}? (\emph{\textbf{Ich war froh, dass …}}) \textquotedblleft{}I hope to visit you\textquotedblright{}? (\emph{\textbf{Ich hoffe, dich zu besuchen.}}) Where does \emph{zu} stand? (\textbf{Right before the infinitive, at the end.})
\end{tcolorbox}
